\documentclass[11pt,letterpaper]{article}
\usepackage[T1]{fontenc}
\usepackage[utf8]{inputenc}
\usepackage{lmodern,geometry,microtype,xcolor,enumitem,hyperref,fancyhdr}
\geometry{margin=0.78in}
\definecolor{ink}{HTML}{17324D}
\definecolor{accent}{HTML}{126A78}
\hypersetup{colorlinks=true,linkcolor=accent,urlcolor=accent,breaklinks=true}
\setlength{\parindent}{0pt}
\setlength{\emergencystretch}{3em}
\setlist[itemize]{leftmargin=*,itemsep=3pt,topsep=3pt}
\pagestyle{fancy}\fancyhf{}\fancyhead[L]{\small Architecture Documentation Suite}\fancyhead[R]{\small v1.0}\fancyfoot[R]{\thepage}\setlength{\headheight}{15pt}
\begin{document}
\thispagestyle{empty}
{\color{accent}\large TECHNICAL REFERENCE \hfill NOVEMBER 2026}\par\vspace{1.1cm}
{\color{ink}\LARGE\bfseries\raggedright Architecture Documentation Review Framework}\par\vspace{0.35cm}
{\Large Question sets, quantitative gates, evidence and remediation}\par\vspace{0.55cm}
\rule{\textwidth}{0.7pt}\par\smallskip
{\small Architecture Documentation Suite | Release 1.0}\par\vspace{0.7cm}
\tableofcontents
\newpage
\section*{1. Review purpose}\addcontentsline{toc}{section}{1. Review purpose}
Review the architecture description for a declared use: development, operational handover, security analysis, change-impact analysis or standards assessment. This is an assessment of documentation fitness; architecture evaluation and implementation testing are separate activities. Adapt the book's chapter 11 six-step procedure to a versioned engineering workflow.\par\smallskip
\section*{2. Procedure and responsibilities}\addcontentsline{toc}{section}{2. Procedure and responsibilities}
\begin{itemize}
\item Establish purpose: sponsor identifies decisions, concerns, success criteria and represented stakeholder roles.
\item Establish subject: freeze source commit, AD version, diagrams, requirements and applicable standards. Record artifact completeness.
\item Build question sets: tailor checks to scope and stakeholder need; define expected answer, evidence, criticality and reviewer.
\item Plan review: appoint lead and independent reviewers, define schedule and evidence format, agree on release gate and exception authority.
\item Conduct review: execute automated checks; independently inspect claims; use scenario walk-throughs to test fitness; record evidence and findings.
\item Analyze results: consolidate disagreements, assign owners and due dates, retest corrections and issue a release disposition.
\end{itemize}
\section*{3. Local acceptance criteria}\addcontentsline{toc}{section}{3. Local acceptance criteria}
The following are proposed suite thresholds, not requirements asserted by the book or ISO. Apply to release scope and record justified N/A decisions. A high-severity defect cannot be offset by unrelated passing checks.\par\smallskip
\begin{itemize}
\item G01 Scope: 100\% of selected views name viewpoint, concern, baseline, purpose and owner. Evidence: view register.
\item G02 Concern coverage: 100\% of high-priority concerns mapped to an applicable view and disposition; no unresolved critical concern. Evidence: trace matrix and review.
\item G03 Notation: 100\% of diagrams have a title, scope and key; all visible elements and edges have names/types and unambiguous meaning. Evidence: diagram inspection.
\item G04 Provenance: 100\% of observed in-scope elements and relationships have resolvable evidence; every inference is marked and every unknown is tracked. Evidence: source register.
\item G05 Correspondence: zero dangling IDs or unexplained conflicting mappings; every deployed runtime element is mapped or explicitly excepted. Evidence: machine checks and human review.
\item G06 Interfaces: every architecture-significant interface covers contract, failure, authorization and data sensitivity or has justified N/A. Evidence: contract review.
\item G07 Behavior: all critical scenarios include success, denial/failure and recovery/termination where applicable. Evidence: traces and state transitions.
\item G08 Rationale: 100\% of significant known decisions have ADR references or explicit unknown disposition; every proposed decision identifies alternatives and tradeoffs.
\item G09 Build: all selected .tex and .puml sources build in the approved toolchain; zero broken references and no unreadable diagrams. Evidence: CI logs plus visual check.
\item G10 Fitness: each represented stakeholder can answer the agreed task questions with evidence; no unresolved blocker. Evidence: signed review records.
\item G11 Currency: AD baseline equals assessed commit or deviation is documented; superseded diagrams and ADRs are marked.
\item G12 Security: each privileged boundary crossing names authenticated identity, policy enforcement, audit and failure behavior or a recorded gap; unresolved critical gaps block security signoff.
\end{itemize}
\section*{4. Question sets by audience}\addcontentsline{toc}{section}{4. Question sets by audience}
Development: where does a feature live? Which modules may depend on it? Which interface changes break consumers? Evidence: package graph, contracts and change-impact trace. Operations: where is the service deployed? Which health indicators, dependencies and recovery transitions apply? Evidence: deployment, runtime and failure traces. Security: where does untrusted data cross authority boundaries? Who can execute a privileged action? Where are audit and redaction enforced? Evidence: security view, policy contracts and denial tests. Architecture governance: are decisions justified, correspondence consistent and concerns covered? Evidence: AD register and review matrix. Standards reviewer: inspect actual licensed requirements and map them independently; a public-summary crosswalk is not sufficient.\par\smallskip
\section*{5. Scoring and dispositions}\addcontentsline{toc}{section}{5. Scoring and dispositions}
For reporting only, coverage equals passing applicable checks divided by all applicable checks. Mark each check PASS, FAIL, NOT-RUN or N/A with evidence and reviewer. Exclude N/A only with a reason approved by review lead. NOT-RUN is not PASS. Release condition: all mandatory applicable gates pass, all critical/high defects resolved or formally accepted by authorized risk owner where permitted, and stakeholder signoff recorded. Do not publish a conformance claim based on this score.\par\smallskip
\section*{6. Finding lifecycle}\addcontentsline{toc}{section}{6. Finding lifecycle}
Use F-ID, gate, severity, observed issue, impact, affected view/element, evidence, owner, remediation, due date, status, retest and exception expiry. Critical: evidence of false security assertion or inability to support essential review; high: missing critical view/contract or conflicting topology; medium: incomplete supporting detail; low: minor clarity defect. Reopen when evidence changes. Example F-01: deployment shows worker access to store but runtime view omits connector; owner verifies configuration and corrects both maps; reviewer retests CR-02/03.\par\smallskip
\section*{7. Automation and limitations}\addcontentsline{toc}{section}{7. Automation and limitations}
Run scripts/validate.py against the JSON manifest to check ID uniqueness, referential integrity, concern mappings and basic provenance fields. Run renderer and TeX build separately. Review actual source evidence, diagram labels, threat assumptions and quality attribute analysis manually. Structural validation is necessary but not sufficient. Re-review changes in trust boundaries, runtime interfaces, topology, responsibilities or architecture drivers.\par\smallskip
\section*{8. Sources}\addcontentsline{toc}{section}{8. Sources}
Clements et al., P.5 and chapter 11 (especially 11.1--11.3); chapter 10 for packet structure. Official C4 notation guidance: \url{https://c4model.com/diagrams/notation}. The local gates, severities and metrics above are original operational extensions.\par\smallskip
\section*{Bibliographic foundation}
Paul Clements et al., \emph{Documenting Software Architectures: Views and Beyond}, 2nd ed., Addison-Wesley/Pearson, copyright 2011, ISBN 978-0-321-55268-6. Section and chapter references refer to the supplied edition. Additional framework sources are identified in the applicable documents.\par
\end{document}
